\scanpage{8}

\deleted{non divisible par $p$} alors \deleted{$Df$} $Df$ a en $x$ un pôle d'ordre au plus $n+1$, et égal à $n+1$ ssi $p\nmid n$.

\deleted{\uline{Corollaire 1} Si $k(s)$ est de car. nulle, alors
\[
\underline{\mathcal L}\otimes\underline{\mathcal L}_n\simeq\underline{\mathcal L}_{n+1}\quad\text{pour tt }n\geq2,
\]
donc
\deleted{$\underline{\mathcal L}_3\otimes\mathcal L\otimes\underline{\mathcal L}_2,\ \mathcal L_4$}
\inserted{$\mathcal L_4\simeq\mathcal L_2\otimes\mathcal L^2$}
\[
\mathcal L_n\simeq\mathcal L^{n-2}\otimes\mathcal L_2,
\]
\deleted{donc} pour $n\geq2$, \deleted{donc \unsure{on a} $\mathcal L_2\simeq$}

\uline{Corollaire 2} Si $k(s)$ est de car $p\neq2$, alors
\[
\deleted{\unsure{\underline{\mathcal L}_3\simeq\mathcal L_2\otimes\mathcal L}},\quad\underline{\mathcal L}_4\simeq\underline{\mathcal L}_3\otimes\mathcal L,\quad\mathcal L_6\simeq\underline{\mathcal L}_5\otimes\underline{\mathcal L}\quad\text{etc}
\]
[$\underline{\mathcal L}_{2n}=\underline{\mathcal L}_{2n-1}\otimes\mathcal L$ si $n\geq\unsure{2}$] enfin
\[
\mathcal L_2^{\unsure{2}}\simeq\mathcal L^{\otimes2}
\]}

\uline{Corollaire} Si $S$ est de car $p>0$ (i.e. $p\underline O_S=0$) et si $n$ est un multiple de la caractéristique, alors $\mathcal L\otimes\underline{\mathcal L}_n\unsure{\longrightarrow|-}\underline{\mathcal L}_{n+1}$ est \uline{nul}.

\deleted{Puisque} Considérons la suite exacte canonique
\diagram{\begin{tikzcd}[column sep=large]
\unsure{0}\arrow[r,dash]&\underline L^{\otimes n}\arrow[r]&\underline L^{\otimes(n+1)}\arrow[r]&(\underline T_{C/S})^{\otimes(n+1)}\otimes_{\underline O_C}\underline O_{S'}\arrow[r]&0
\end{tikzcd}}
on en conclut la suite exacte
\diagram{\begin{tikzcd}[column sep=large]
0\arrow[r]&\mathcal E_n\arrow[r,dash]&\mathcal E_{n+1}\arrow[r]&\mathcal L^{\otimes(n+1)}\arrow[r]&\unsure{R^1}(f,\underline L^{\otimes n})\arrow[d,equal]\\
&&&&0\ \text{si }n\geq1
\end{tikzcd}}
d'où
\[
\boxed{\mathcal L_{\unsure{n+1|n}}\simeq\mathcal L^{\otimes(n+1)}\qquad n\geq2}
\]
\margin{\unsure{intervertir} avec la prop. 2}
et la multiplication dans l'\unsure{A} graduée de prop. 2 n'est autre que la multiplication triviale pour cette forme des $\mathcal L_n$, $n\geq2$.

\mspage{7}
